\section{Discussion}
The version-2 logs reinforce the main motivation of SynFoC: U-Net and MedSAM contribute differently across mixed-domain datasets, so a joint pipeline is more reliable than assuming one branch is always stronger. On Prostate and M\&Ms, the reproduced U-Net branch reaches the best average DSC, while on Fundus and BUSI, the MedSAM branch gives the better overall averages. This pattern is consistent with the idea that the stronger pseudo-label source depends on both the dataset characteristics and the training stage.

The qualitative figures also support the quantitative results. Across the four datasets, the reproduced masks are generally more complete on target structures when the two branches agree, while disagreement regions remain the main source of boundary errors. This makes the self-mutual confidence rule and the consensus-divergence consistency objective useful not only as optimization details, but also as practical mechanisms for handling cross-domain instability.

The current report still has several limitations. First, the conclusions are drawn from the completed logged runs with 40 labeled samples from domain 1, so comparisons with published baselines should be interpreted as contextual rather than strictly matched settings. Second, the analysis is based on the best recorded checkpoints in the available logs, not on repeated runs with statistical significance testing. Third, this work is a reproduction study, so it does not yet propose a new architectural extension beyond the released SynFoC pipeline. Future work should validate the stability of the reproduced results across multiple seeds and then explore stronger uncertainty estimation, more robust prompt generation, or lighter adaptation strategies for the foundation branch.
